
%\needspace{4\baselineskip}
\color{uniblue}{\subsection{Токовая отсечка (ТО)}}
\color{black}

\begin{enumerate}

\item
Токовая отсечка применяется для защиты обмоток трансформатора от повреждений внутри трансформатора, сопровождающихся большими токами. Для трансформаторов до 35 кВ мощностью менее 6,3 МВА в некоторых случаях допускается применять токовую отсечку вместо дифференциальной токовой защиты трансформатора. В устройстве токовая отсечка представлена функцией «ТО».

\item
Ввод функции в работу на этапе параметрирования устройства осуществляется программным переключателем SGF1 «Ввод\_функции». Информация о введенной функции отображается выходным сигналом «ТО: Ввод».
\item
Функция может быть выведена из работы оперативно путем активации сигнала «ОВ ТО», что характеризуется активным состоянием сигнала «ТО: Оперативный вывод» на выходе функции. 
\item
При выборе положения «Треугольник» программного переключателя SGF2 «Сборка\_ток\_цепей» происходит исключение токов нулевой последовательности из расчета рабочих токов. 
\item
Функция контролирует повышение любого фазного тока ($I_{A\ \textsl{мгн}}, I_{B\ \textsl{мгн}}, I_{C\ \textsl{мгн}}$) выше заданной уставки, при превышении заданной уставки по току через выдержку времени таймера T1 (уставка «Тср») выдается сигнал на отключение трансформатора. При активном состоянии входного сигнала «ТО на сигн» действие токовой отсечки переводится на сигнал.

\item
Алгоритм работы функции выполнен в соответствии с рисунком \ref{to:img1}. В таблице \ref{to:tbl1} приведены параметры, необходимые для настройки функции.

\vspace{1mm}
\begin{figure}[H]
\centering
\includegraphics[width=1\textwidth,height=1\textheight,keepaspectratio]{\fbpath/29. ТО 35/_latex/img/to.pdf}
\captionof{figure}{Функциональная схема логики ТО}
\label{to:img1}
\end{figure}

\end{enumerate}
\FloatBarrier % Все float'ы ДО этой строки будут размещены здесь.